This element simulates a beam-driven monopole mode cavity using the fundamental theorem of beam loading and phasor rotation.
In addition, a generator-driven field may be included using a feedback system \cite{Berenc-IPAC15-MOPMA006}.

Note on phase conventions: the phase convention for the \verb|PHASE_SETPOINT|  parameter of \verb|RFMODE| is the
same as for the \verb|PHASE| parameter of \verb|RFCA|. However, in the output files from \verb|RFMODE|, i.e., the
files requested with the \verb|RECORD| and \verb|FEEDBACK_RECORD| parameters, a different convention is used, which 
differs by $-90$ degrees from the \verb|PHASE_SETPOINT|  parameter. 

The feedback implementation uses either amplitude and phase feedback or else in-phase and quadrature feedback.
Figure \ref{fig:rfFeedbackModel} shows the model used for the feedback system.
More information is available in \cite{Berenc-IPAC15-MOPMA006}.

Rf feedback is active when a non-zero value is given for \verb|DRIVE_FREQUENCY| and when either
\verb|AMPLITUDE_FILTER| and \verb|PHASE_FILTER| or else
\verb|IN_PHASE_FILTER| and \verb|QUADRATURE_FILTER| are given.
These parameters name SDDS files that define filters in the $z$ domain 
\begin{equation}
H(z) = \frac{\sum\limits_{i=0}^m b_i z^{-i}}{\sum\limits_{i=0}^r a_i z^{-i}}
\end{equation}
which translates into the discrete time domain as
\begin{equation}
y[n ] = -\frac{1}{a_0}\sum\limits_{i=1}^r a_i y_{r-i} + \frac{1}{a_0}\sum\limits_{i=0}^m b_i x_{n-i},
\end{equation}
where $y[n]$ is the filter output (e.g., additional generator current) and $x[n]$ is the filter input (e.g., 
the voltage error signal). 
The filter files must each contain two columns:
\begin{itemize}
\item \verb|denominator| --- The $a_i$ filter coefficients. These should generally sum to zero or nearly zero.
\item \verb|numerator| --- The $b_i$ filter coefficients. These are generally very small.
\end{itemize}
Each file can consist of up to four pages, with each page representing one filter in a parallel filter bank
formed from all pages. The output of the four stages is added to obtain the control signal.

Optionally, each page may define a floating-point parameter named \verb|gain|. When present, the output of that
page's filter stage is multiplied by its \verb|gain| value before being passed on---i.e. before being summed into
the parallel bank (or, where the filter is applied as a series cascade, before being fed to the next stage). The
unscaled output is retained internally as the stage's feedback history, so the stage's own transfer function
$H(z)$ is unchanged and only the signal passed on is scaled. When the \verb|gain| parameter is absent it defaults
to unity, so filter files that do not define it behave exactly as before.
The number of rows in the page is given by the larger of r or m; in the case of multiple pages,
the number of rows in each page is given by the largest r or m from the highest order filter amongst the pages, by
simply stuffing the rows of the other pages with zeros, i.e., rows representing orders of $z$ down to $z^{-l}$ where $l$ is
the largest $r$ or $m$ of the filter pages.

For example, a digital low-pass filter with DC gain of $k$ can be specified as
\begin{equation}
y[1] = e^{-pT} y[0] + k (1 - e^{-pT}) x[0],
\end{equation}
where $p$ is the frequency of the pole and $T$ is the sample interval.

The feedback loop reads the cavity state and acts on the generator at a fixed interval (in buckets) of
\verb|UPDATE_INTERVAL|. 
The timing of this activity is aligned to the arrival time of the first bunch in the \verb|RFMODE| element.
By default (\verb|READ_OFFSET|=0), the timing is such that the state is read just before the next arrival of
that bunch; in particular, it is 180 degrees ahead of that arrival.
If bunches are equally spaced by, say $N_b$ buckets, the \verb|UPDATE_INTERVAL| parameter should ideally be
$m N_b$, where $m>0$ is an integer.
This ensures that the state is read at a fixed timing relative to the bunches.

The rf feedback feature makes use of the voltage amplitude measured when there is no bunch present.
The \verb|RECORD| file shows the voltage seen by the beam, computed by averaging over the voltage for
each particle.
These may deviate by values from a few percent to of order ten percent, depending on the loss factor for the
cavity and the number of bunchess; this is caused by the fact that the rate at which an intense bunch removes
energy from the cavity will typically, albeit briefly, exceed the power from the generator.
To reduce the impact of this effect, one may use the \verb|ADJUSTMENT_FRACTION|, \verb|ADJUSTMENT_START|, and \verb|ADJUSTMENT_INTERVAL|
parameters to modify the voltage setpoint.
If \verb|ADJUSTMENT_FRACTION| is non-zero, then for every \verb|ADJUSTMENT_INTERVAL|$^{th}$ pass after the 
\verb|ADJUSTMENT_START|$^{th}$ pass, the voltage setpoint will be adjusted based on a comparison of the bunch-averaged
voltage to the user's setpoint.
E.g., if the bunch-averaged voltage is 100 V too low and \verb|ADJUSTMENT_FRACTION| is 0.1, the voltage setpoint will
be raised by 10 V.
Users should note that if \verb|ADJUSTMENT_FRACTION| is too large or \verb|ADJUSTMENT_INTERVAL| is too small, the system
may be unstable.

\begin{figure}[htb]
\center
\includegraphics{rfFeedbackModel}
\caption{Rf feedback model used by the {\tt RFMODE} element. Courtesy T. Berenc, Argonne National Laboratory.}
\label{fig:rfFeedbackModel}
\end{figure}

Normally, the field dumped in the cavity by one particle affects trailing particles in the same turn.
However, if one is also using a \verb|WAKE| or \verb|ZLONGIT| element to simulate the short-range wake of the cavity, this would be double-counting.
In that case, one can use \verb|LONG_RANGE_ONLY=1| to suppress the same-turn effects of the \verb|RFMODE| element.

The \verb|FREQ_WAVEFORM| and \verb|Q_WAVEFORM| parameters scale the mode's resonant frequency and quality
factor as functions of time. The time column of each file is interpreted as laboratory time, i.e., the full
elapsed time since the start of tracking, including any time removed by an \verb|RFCA| or \verb|RFCW|
element with \verb|CHANGE_T|=1.

Users should be aware of a limitation when either waveform is combined with \verb|CHANGE_T|=1. Reduced-time
tracking removes close to a revolution period from each particle's time coordinate every turn, and the phasor
rotation must credit the removed time with the phase it would have accumulated. That correction is exact only
while the mode's angular frequency is constant, which a frequency waveform (and, far more weakly, a $Q$
waveform) violates. The resulting error in the beam-induced voltage grows roughly in proportion to the
fractional change in frequency per turn, so it is negligible for a slow ramp and significant for a fast one.
{\tt elegant} issues a warning when it detects this combination. If the ramp is appreciable, either omit
\verb|CHANGE_T|=1 or drop the waveform.

Two output files are available: the \verb|RECORD| file includes bunch-by-bunch data on the beam-induced fields and the total cavity fields.
The \verb|FEEDBACK_RECORD| file includes tick-by-tick data from the feedback system simulation; {\em writing this file this can significantly degrade performance.}

NB: when \verb|BUNCHED_BEAM_MODE| is set to a value other than 1, in order to obtain the effect of several bunches while tracking
only one bunch, the total charge set with the \verb|TOTAL| parameter of the \verb|CHARGE| element should equal the charge in
a single bunch, not the entire beam. However, when \verb|BUNCHED_BEAM_MODE|=1 (allowing an indeterminant number of bunches to be
actually present), then \verb|TOTAL| should be the total for all bunches together.
